% Drawing 147/05-4 -- folder 147, batch 05, archive page 97 (PDF page 98).
% Transcription: le graphe $G'$ obtenu en coupant aux croix : le sommet trivalent, l'arête coupée devenue deux boucles aplaties, la boucle coupée devenue deux boucles aplaties issues du troisième sommet.
% Kind: graph G' obtained by cutting at the crosses. Method 2 (schematic redraw from the facsimile). Status: ok.
% Notes: G': the trivalent vertex with the vertex above carrying two flattened loops and the vertex below left; the cut edge has become two flattened loops (one from the trivalent vertex, one from the second vertex), and the cut loop two flattened loops from the third vertex (with its third free segment).
% Coordinates: 1 unit = 1 mm = 1/100 of the crop width (bbox in page fractions: [0.235, 0.49, 0.531, 0.563]).
% Redrawn by Claude (claude-opus-5-5), issue #73, 2026-11-07; not checked by a human; not compiled (no TeX in the session).
% Rights: reproduces part of the work -- behind RENDER_TRANSCRIPTS, never published by itself.
\documentclass[tikz,border=3pt]{standalone}
\usepackage{amsmath,amssymb}
\usetikzlibrary{patterns,arrows.meta}
\tikzset{pencil/.style={black!45}, struck/.style={black!35}, mark/.style={circle,draw,inner sep=1.3pt}, pt/.style={circle,fill,inner sep=0.9pt}, opt/.style={circle,draw,fill=white,inner sep=1.1pt}}
\begin{document}
\begin{tikzpicture}[x=1mm,y=1mm,line cap=round,line join=round,font=\scriptsize]
\node[pt] at (16.3,-11.3) {};
\draw (16.3,-11.3) -- (15,-3.8);
\draw (16.3,-11.3) -- (21.3,-3.8);
\draw (16.3,-11.3) -- (20,-22.5);
\node[pt] at (20,-22.5) {};
\draw (20,-22.5) -- (16.3,-30);
\node[pt] at (16.3,-30) {};
\draw (20,-22.5) -- (32.5,-22.5);
\draw (45,-22.5) -- (57.5,-21.3);
\node[pt] at (57.5,-21.3) {};
\draw (57.5,-21.3) -- (68.8,-11.3);
\node[pt] at (68.8,-11.3) {};
\draw (68.8,-11.3) -- (81.3,-11.3);
\draw (68.8,-11.3) -- (70,-21.3);
\draw (68.8,-11.3) -- (78.8,-18.8);
\end{tikzpicture}
\end{document}
